% Source: https://artofproblemsolving.com/community/c3h1212777_mc_problems
% Stable identity: aops-topic-1212777-q01
The fingers form a row of $5$ positions; we count permutations of $1,\dots,5$ in which no two consecutive integers are adjacent. Look at where $3$ is: its neighbours can only be $1$ and $5$.
\begin{itemize}
\item $3$ in the middle: $\_,1,3,5,\_$ or $\_,5,3,1,\_$. The ends get $2$ and $4$; $2$ cannot be next to $1$ and $4$ cannot be next to $5$, so $4,1,3,5,2$ and $2,5,3,1,4$: $2$ ways.
\item $3$ in position $2$: then $\{1,5\}$ fill positions $1$ and $3$. If position $3$ is $5$, the last two are $2,4$ in an order with $4$ not next to $5$: $1,3,5,2,4$. If position $3$ is $1$: $5,3,1,4,2$. So $2$ ways, and by reversal $2$ ways with $3$ in position $4$: $4$ ways.
\item $3$ at an end, say position $1$: position $2$ is $1$ or $5$. With $3,1$: the rest $\{2,4,5\}$ must avoid $1$--$2$, $4$--$5$: $3,1,4,2,5$ and $3,1,5,2,4$. With $3,5$: by the symmetry $k\mapsto6-k$, also $2$ ways. So $4$ ways, and $4$ more with $3$ at the other end: $8$ ways.
\end{itemize}
Total $2+4+8=\boxed{14}$.
